\documentclass[11pt]{article}

\usepackage[T1]{fontenc}
\usepackage[margin=1in]{geometry}
\usepackage{amsmath,amssymb,amsthm}
\usepackage{url}

\newcommand{\R}{\mathbb{R}}
\newcommand{\Prob}{\mathbb{P}}
\newcommand{\Frob}[1]{\left\lVert #1\right\rVert_F}
\newcommand{\Norm}[1]{\left\lVert #1\right\rVert_2}
\newcommand{\Atil}{\widetilde A}
\newcommand{\Util}{\widetilde U}
\newcommand{\Ahat}{\widehat A}
\newcommand{\Ghat}{\widehat G}
\newcommand{\I}{I}
\newcommand{\eps}{\varepsilon}
\newcommand{\One}{\mathbf 1}

\newtheorem{theorem}{Theorem}
\newtheorem*{baseline}{Exact baseline}
\theoremstyle{remark}
\newtheorem*{remark}{Remark}

\title{Algorithm 2: Floating-Point Changes to Equations (5) and (7)}
\author{}
\date{}

\begin{document}
\sloppy

\maketitle

\section*{Notation and Arithmetic Model}

Let \(A\in\R^{m\times n}\).  The \(i\)-th row of \(A\) is \(A_{i*}\), and
\[
  D(A)=\Frob{A}^2=\sum_{i=1}^m \Norm{A_{i*}}^2 .
\]
Assume \(D(A)>0\).  Algorithm 2 samples indices
\[
  I_1,\ldots,I_s\in\{1,\ldots,m\}
\]
independently from the exact row-norm probability law
\[
  p_i=\frac{\Norm{A_{i*}}^2}{D(A)}.
\]
The probability symbol \(\Prob[\cdot]\) is always taken with respect to this
exact product sampling law.  This is the row-sampling law from Algorithm 2 in
Drineas and Mahoney's RandNLA survey \cite{DrineasMahoney2016RandNLA}.

The exact sampled sketch is \(\Atil\in\R^{s\times n}\), with entries
\[
  \Atil_{tj}=\frac{A_{I_tj}}{\sqrt{s\,p_{I_t}}}.
\]
The exact Gram matrix of a matrix \(B\) is \(B^TB\).  Thus the exact sampled
Gram matrix is \(\Atil^T\Atil\), and the exact reference Gram matrix is
\(A^TA\).

The floating-point unit roundoff is \(u\ge0\).  When a Gram entry is computed
by a length-\(s\) floating-point dot product, \(\gamma_s\) denotes the usual
valid dot-product error factor.  The probabilities \(p_i\) and the sampling law
are exact by convention, but the non-probability denominator
\[
  d_i=\sqrt{s\,p_i}
\]
is computed.  The concrete denominator routine formalized here is
\[
  \widehat d_i=\operatorname{flsqrt}(\operatorname{flmul}(s,p_i)).
\]
It satisfies, for every positive-probability row,
\[
  |\widehat d_i-d_i|
  \le
  d_i\,\alpha(u),
  \qquad
  \alpha(u)=\sqrt{1+u}\,u+u,
  \qquad
  \widehat d_i\ne0.
\]
Define the support relative denominator budget and the resulting effective
row-scaling relative error by
\[
  \rho_A
  =
  \sum_{i:\,p_i>0}
    \frac{d_i\,\alpha(u)}{|\widehat d_i|},
  \qquad
  \eta_A
  =
  u+(1+u)\rho_A .
\]
The deterministic perturbation budgets used below are
\[
  \tau_{\rm scale}^{\rm comp}(A)
  =
  n(2\eta_A+\eta_A^2)D(A),
\]
\[
  \tau_{\rm dot}^{\rm comp}(A;s)
  =
  n\gamma_s(1+\eta_A)^2D(A),
\]
and
\[
  \tau_{\rm full}^{\rm comp}(A;s)
  =
  \tau_{\rm scale}^{\rm comp}(A)+\tau_{\rm dot}^{\rm comp}(A;s).
\]
For comparison, if \(d_i\) is supplied exactly, then \(\rho_A=0\),
\(\eta_A=u\), and these reduce to the older exact-denominator budgets.

\section*{What Is Exact and What Is Computed}

\begin{itemize}
  \item Exact arithmetic inputs: the matrix \(A\), the probabilities \(p_i\),
  the random samples \(I_t\), the exact sketch \(\Atil\), and the exact
  reference Gram matrix \(A^TA\).
  \item Computed denominator table: \(\widehat d_i\) is computed by
  \(\operatorname{flsqrt}(\operatorname{flmul}(s,p_i))\).  This is a
  floating-point non-probability quantity and its error contributes through
  \(\rho_A\) and \(\eta_A\).
  \item Floating-point row-scaling output:
  \(\Ahat_{\rm comp}\in\R^{s\times n}\) denotes the sketch whose entries are
  produced as \(\operatorname{fldiv}(A_{I_tj},\widehat d_{I_t})\).
  \item Exact arithmetic after a floating-point sketch: the matrix
  \(\Ahat_{\rm comp}^T\Ahat_{\rm comp}\) is the exact mathematical Gram matrix
  of the already rounded sketch.  No dot-product kernel is charged in this
  model.
  \item Floating-point computed Gram output: \(\Ghat_{\rm dot}(A)\) denotes the
  Gram matrix whose entries are computed by floating-point dot products applied
  to the computed-denominator rounded sketch \(\Ahat_{\rm comp}\).  This model
  charges denominator computation, rounded row scaling, and dot products through
  \(\tau_{\rm full}^{\rm comp}(A;s)\).
  \item Equation (7) exact analysis basis: the matrix \(U\) is the exact
  orthonormal-column basis used to define the leverage probabilities, the exact
  sampling law, and the reference matrix \(\I_n\).
  \item Equation (7) computed basis table: the final stored-basis theorem below
  uses a rounded table \(\widehat U^{(\mu)}\) in the sampled sketch and charges
  its entrywise storage error, the computed denominator table, rounded row
  divisions, and rounded Gram dot products.  The three closed storage modes are
  \(\mu\in\{\mathrm{mul1},\mathrm{add0},\mathrm{sub0}\}\).
  \item Equation (7) actual-input endpoint: when an exact analysis
  factorization \(A=UC\) is supplied with \(U^TU=I\), the theorem below samples
  rows of the actual input \(A\) using the exact leverage law defined by \(U\).
  The computed non-probability quantities are the concrete denominator table,
  the rounded row divisions on entries of \(A\), and the rounded Gram dot
  products.  The factors \(U\) and \(C\) are exact analysis/probability
  witnesses in that theorem, not computed algorithm outputs.
  \item Equation (7) basis generation from a data matrix: the theorem below
  does not claim that a QR, SVD, or rank-revealing routine has computed \(U\)
  from a matrix \(A\).  Such a routine must provide its own computed-basis
  theorem before it can be composed with this stored-basis endpoint.
\end{itemize}

\section*{Equation (5)}

\begin{baseline}[Equation (5), exact arithmetic]
This is equation (5) in the RandNLA survey
\cite{DrineasMahoney2016RandNLA}.
For \(s>0\) and \(\eps>0\),
\[
  \Prob\!\left[
    \Frob{\Atil^T\Atil-A^TA}
    \le \eps D(A)
  \right]
  \ge
  1-\frac{1}{s\eps^2}.
\]
\end{baseline}

\begin{theorem}[Equation (5), computed denominators and rounded row scaling]
For \(s>0\) and \(\eps>0\),
\[
  \Prob\!\left[
    \Frob{\Ahat_{\rm comp}^T\Ahat_{\rm comp}-A^TA}
    \le
    \eps D(A)+\tau_{\rm scale}^{\rm comp}(A)
  \right]
  \ge
  1-\frac{1}{s\eps^2}.
\]
Thus the exact equation (5) sampling probability is unchanged, while the
threshold is enlarged by the proved deterministic denominator-plus-division
budget.
\end{theorem}

\begin{theorem}[Equation (5), fully computed denominators and Gram]
Assume the length-\(s\) dot-product factor \(\gamma_s\) is valid.  For \(s>0\)
and \(\eps>0\),
\[
  \Prob\!\left[
    \Frob{\Ghat_{\rm dot}(A)-A^TA}
    \le
    \eps D(A)+\tau_{\rm full}^{\rm comp}(A;s)
  \right]
  \ge
  1-\frac{1}{s\eps^2}.
\]
The additive floating-point part has order
\[
  \tau_{\rm full}^{\rm comp}(A;s)
  =
  \Theta\!\left(
    nD(A)\bigl(\eta_A+\eta_A^2+\gamma_s(1+\eta_A)^2\bigr)
  \right).
\]
In the small-roundoff regime where \(\eta_A\to0\) and \(\gamma_s\to0\),
\[
  \tau_{\rm full}^{\rm comp}(A;s)
  =
  \Theta\!\left(nD(A)(\eta_A+\gamma_s)\right).
\]
For the concrete denominator routine above, \(\alpha(u)=\Theta(u)\), and
\(\eta_A\) is determined by \(u\), the exact probabilities, and the computed
denominators \(\widehat d_i\).  For fixed \(A,m,n,s\) and nonzero computed
positive-probability denominators, the additive FP budget tends to zero as
\(u,\gamma_s\to0\), so the result is non-vacuous.
\end{theorem}

\section*{Equation (7)}

Now let \(U\in\R^{m\times n}\) have orthonormal columns:
\[
  U^TU=\I_n.
\]
Assume \(n>1\).  Algorithm 2 samples rows from the exact leverage-score law
\[
  p_i^U=\frac{\Norm{U_{i*}}^2}{n}.
\]
Let \(\Util\in\R^{s\times n}\) be the exact leverage-sampled sketch.  For
symmetric matrices \(X,Y\in\R^{n\times n}\), write \(X\preceq Y\) when
\[
  x^T X x\le x^T Y x
  \qquad\hbox{for every }x\in\R^n.
\]

\begin{baseline}[Equation (7), exact Bennett form]
This is the leverage-score subspace-embedding bound corresponding to equation
(7) in the RandNLA survey \cite{DrineasMahoney2016RandNLA}.  The underlying
leverage-sampling theorem appears, for example, in
Drineas--Mahoney--Muthukrishnan--Sarlos \cite{DrineasMahoneyMuthukrishnanSarlos2011}
and is stated here through the rank-one covariance concentration route of
Oliveira \cite{Oliveira2010} and Tropp-style matrix concentration
\cite{Tropp2012}.
Assume \(s>0\), \(\eps>0\), \(\delta>0\), and
\[
  \log\frac{2n}{\delta}
  \le
  s\,\frac{\eps^2}{2(n-1)+(2/3)n\eps},
  \qquad
  \log\frac{2n}{\delta}
  \le
  s\,\frac{\eps^2}{2(n-1)+(2/3)\eps}.
\]
Then
\[
  \Prob\!\left[
    \Util^T\Util-\I_n\preceq \eps\I_n
    \ \hbox{and}\
    \I_n-\Util^T\Util\preceq \eps\I_n
  \right]\ge 1-\delta .
\]
\end{baseline}

\begin{theorem}[Equation (7), computed-denominator floating-point Bennett form]
Assume the hypotheses and sample-budget inequalities in the exact Bennett
baseline, and assume the length-\(s\) dot-product factor \(\gamma_s\) is valid.
Let \(\Ghat_{\rm dot}(U)\) be the floating-point computed Gram matrix formed
from the leverage-sampled sketch using the concrete computed denominator
\(\widehat d_i^U=\operatorname{flsqrt}(\operatorname{flmul}(s,p_i^U))\),
rounded row scaling, and rounded length-\(s\) Gram dot products.  Then
\[
  \Prob\!\left[
    \Ghat_{\rm dot}(U)-\I_n
      \preceq
      \bigl(\eps+\tau_{\rm full}^{\rm comp}(U;s)\bigr)\I_n
    \ \hbox{and}\
    \I_n-\Ghat_{\rm dot}(U)
      \preceq
      \bigl(\eps+\tau_{\rm full}^{\rm comp}(U;s)\bigr)\I_n
  \right]\ge 1-\delta .
\]
Since \(U^TU=\I_n\), \(D(U)=\Frob{U}^2=n\), and therefore
\[
  \tau_{\rm full}^{\rm comp}(U;s)
  =
  n^2\bigl(2\eta_U+\eta_U^2+\gamma_s(1+\eta_U)^2\bigr).
\]
Its small-roundoff order is
\[
  \tau_{\rm full}^{\rm comp}(U;s)
  =
  \Theta\!\left(n^2(\eta_U+\gamma_s)\right),
  \qquad
  \eta_U=u+(1+u)
  \sum_{i:\,p_i^U>0}
    \frac{\sqrt{s p_i^U}\,\alpha(u)}
         {|\operatorname{flsqrt}(\operatorname{flmul}(s,p_i^U))|}.
\]
\end{theorem}

\subsection*{Actual Input Matrix for Equation (7)}

This subsection states the end-to-end row-sampling conclusion for the actual
matrix whose leverage scores are analyzed.  To avoid overloading dimensions,
let \(U\in\R^{m\times r}\) have orthonormal columns, let
\(C\in\R^{r\times q}\), and define
\[
  A=UC\in\R^{m\times q}.
\]
The exact leverage probabilities are
\[
  p_i^U=\frac{\Norm{U_{i*}}^2}{r}.
\]
The factors \(U\) and \(C\) are exact analysis/probability witnesses here:
the theorem does not assert that a QR, SVD, or rank-revealing routine has
computed them.  The implemented sampled sketch uses rows of \(A\), not rows of
\(U\).  Its exact row \(t\) is
\[
  \widetilde A_{U,t*}
  =
  \frac{A_{I_t*}}{\sqrt{s p_{I_t}^U}},
  \qquad I_t\sim p^U .
\]
Assume \(r>1\), \(s>0\), \(\eps>0\), \(\delta>0\), and
\[
  \log\frac{2r}{\delta}
  \le
  s\,\frac{\eps^2}{2(r-1)+(2/3)r\eps},
  \qquad
  \log\frac{2r}{\delta}
  \le
  s\,\frac{\eps^2}{2(r-1)+(2/3)\eps}.
\]
Then exact arithmetic gives the relative Gram statement
\[
  \Prob\!\left[
    \widetilde A_U^T\widetilde A_U-A^TA
      \preceq \eps A^TA
    \ \hbox{and}\
    A^TA-\widetilde A_U^T\widetilde A_U
      \preceq \eps A^TA
  \right]\ge 1-\delta .
\]

For floating-point arithmetic, the concrete leverage denominator is
\[
  \widehat d_i^U
  =
  \operatorname{flsqrt}(\operatorname{flmul}(s,p_i^U)),
\]
and the effective row-scaling relative error is
\[
  \eta_U^A
  =
  u+(1+u)
  \sum_{i:\,p_i^U>0}
  \frac{\sqrt{s p_i^U}\,(\sqrt{1+u}\,u+u)}
       {|\operatorname{flsqrt}(\operatorname{flmul}(s,p_i^U))|}.
\]
For each actual-input column \(j\), define
\[
  L_j^A
  =
  \sum_{i:\,p_i^U>0}\frac{|A_{ij}|}{\sqrt{p_i^U}}.
\]
The fully computed actual-input radius is
\[
  T_A^{\rm fp}
  =
  \left(
    \gamma_s(1+\eta_U^A)^2+2\eta_U^A+(\eta_U^A)^2
  \right)
  \sum_{j=1}^q (L_j^A)^2 .
\]
With all local shorthand expanded, this is
\[
\begin{aligned}
T_A^{\rm fp}
&=
\Biggl[
\gamma_s
\left(
1+u+(1+u)
  \sum_{i:\,\Norm{U_{i*}}^2/r>0}
  \frac{\sqrt{s\,\Norm{U_{i*}}^2/r}\,(\sqrt{1+u}\,u+u)}
       {|\operatorname{flsqrt}(\operatorname{flmul}(s,\Norm{U_{i*}}^2/r))|}
\right)^2 \\
&\qquad
+2\left(
u+(1+u)
  \sum_{i:\,\Norm{U_{i*}}^2/r>0}
  \frac{\sqrt{s\,\Norm{U_{i*}}^2/r}\,(\sqrt{1+u}\,u+u)}
       {|\operatorname{flsqrt}(\operatorname{flmul}(s,\Norm{U_{i*}}^2/r))|}
\right) \\
&\qquad
+\left(
u+(1+u)
  \sum_{i:\,\Norm{U_{i*}}^2/r>0}
  \frac{\sqrt{s\,\Norm{U_{i*}}^2/r}\,(\sqrt{1+u}\,u+u)}
       {|\operatorname{flsqrt}(\operatorname{flmul}(s,\Norm{U_{i*}}^2/r))|}
\right)^2
\Biggr] \\
&\qquad\qquad\cdot
\sum_{j=1}^q
\left(
  \sum_{i:\,\Norm{U_{i*}}^2/r>0}
  \frac{\left|\sum_{\ell=1}^r U_{i\ell}C_{\ell j}\right|}
       {\sqrt{\Norm{U_{i*}}^2/r}}
\right)^2 .
\end{aligned}
\]

\begin{theorem}[Equation (7), actual input matrix and computed Gram]
Under the hypotheses above, assume the length-\(s\) dot-product factor
\(\gamma_s\) is valid.  Let \(\Ghat_{\rm dot}(A;U)\) be the floating-point
Gram computed by sampling exact rows of \(A=UC\) from the exact leverage law
defined by \(U\), computing
\(\widehat d_i^U=\operatorname{flsqrt}(\operatorname{flmul}(s,p_i^U))\),
rounding the sampled-row divisions, and computing the Gram entries by
floating-point dot products.  Then
\[
  \Prob\!\left[
    \Ghat_{\rm dot}(A;U)-A^TA
      \preceq
      \eps A^TA+T_A^{\rm fp}\I_q
    \ \hbox{and}\
    A^TA-\Ghat_{\rm dot}(A;U)
      \preceq
      \eps A^TA+T_A^{\rm fp}\I_q
  \right]\ge 1-\delta .
\]
\end{theorem}

\begin{remark}
Let \(k_U=|\{i:p_i^U>0\}|\).  In the fixed-problem small-roundoff regime where
\(U,C,m,r,q,s\) are fixed, positive-probability computed denominators remain
nonzero and converge to their exact values, and \(u,\gamma_s\to0\),
\[
  \eta_U^A=\Theta(u(1+k_U)).
\]
Consequently
\[
  T_A^{\rm fp}
  =
  \Theta\!\left(
    \bigl(\gamma_s+\eta_U^A\bigr)
    \sum_{j=1}^q (L_j^A)^2
  \right)
  =
  \Theta\!\left(
    \left(\gamma_s+u(1+k_U)\right)
    \sum_{j=1}^q
    \left(
      \sum_{i:\,p_i^U>0}
      \frac{|A_{ij}|}{\sqrt{p_i^U}}
    \right)^2
  \right).
\]
The radius tends to zero in this regime, so the computed actual-input theorem
is non-vacuous.  If a concrete implementation first computes \(U\) or \(C\),
that QR/SVD/rank-revealing routine must be certified separately before this
exact-analysis-witness theorem is claimed for that generation path.
\end{remark}

\subsection*{Stored Basis Table for Equation (7)}

The basis-sampled computed-denominator theorem above charges the computed
denominator, row divisions, and dot products while using the exact entries of
\(U\) in the sampled sketch.  The next result is the implementation-facing
stored-basis endpoint for the modeled rounded-copy paths.  The exact basis
\(U\) still defines the probability law and the reference identity matrix, but
the computed sketch uses a stored table.

Let
\[
  \mathcal M=\{\mathrm{mul1},\mathrm{add0},\mathrm{sub0}\}.
\]
For \(\mu\in\mathcal M\), define the stored basis table by
\[
  \widehat U^{(\mathrm{mul1})}_{ij}
    =\operatorname{flmul}(U_{ij},1),
  \qquad
  \widehat U^{(\mathrm{add0})}_{ij}
    =\operatorname{fladd}(U_{ij},0),
  \qquad
  \widehat U^{(\mathrm{sub0})}_{ij}
    =\operatorname{flsub}(U_{ij},0).
\]
For each of these three storage modes,
\[
  |\widehat U^{(\mu)}_{ij}-U_{ij}|\le u|U_{ij}|.
\]
The leverage denominator is still computed by
\[
  d_i^U=\sqrt{s\,p_i^U},
  \qquad
  \widehat d_i^U
  =
  \operatorname{flsqrt}(\operatorname{flmul}(s,p_i^U)).
\]
On positive-probability rows, the concrete denominator routine proves
\[
  |\widehat d_i^U-d_i^U|\le d_i^U\,\alpha(u),
  \qquad
  \widehat d_i^U\ne0.
\]
Define the denominator relative budget
\[
  \rho_U
  =
  \sum_{i:\,p_i^U>0}
  \frac{\sqrt{s p_i^U}\,(\sqrt{1+u}\,u+u)}
       {|\operatorname{flsqrt}(\operatorname{flmul}(s,p_i^U))|}.
\]
The total relative sampled-entry budget from denominator formation, rounded
division, and one rounded-copy storage of the basis table is
\[
  \lambda_U
  =
  2u+u^2+(1+u)^2\rho_U .
\]
For each column \(j\), define the exact positive-support column magnitude
\[
  B_j^U
  =
  \sum_{i:\,p_i^U>0}\frac{|U_{ij}|}{\sqrt{p_i^U}}.
\]
The stored-basis sampled-Gram perturbation radius is
\[
  T_U^{\rm fp}
  =
  \left(
    \gamma_s(1+\lambda_U)^2+2\lambda_U+\lambda_U^2
  \right)
  \sum_{j=1}^n (B_j^U)^2 .
\]
Equivalently, with every shorthand replaced,
\[
\begin{aligned}
T_U^{\rm fp}
&=
\Biggl[
\gamma_s
\left(
  1+
  2u+u^2
  +(1+u)^2
  \sum_{i:\,p_i^U>0}
  \frac{\sqrt{s p_i^U}\,(\sqrt{1+u}\,u+u)}
       {|\operatorname{flsqrt}(\operatorname{flmul}(s,p_i^U))|}
\right)^2 \\
&\qquad
+2\left(
  2u+u^2
  +(1+u)^2
  \sum_{i:\,p_i^U>0}
  \frac{\sqrt{s p_i^U}\,(\sqrt{1+u}\,u+u)}
       {|\operatorname{flsqrt}(\operatorname{flmul}(s,p_i^U))|}
\right) \\
&\qquad
+\left(
  2u+u^2
  +(1+u)^2
  \sum_{i:\,p_i^U>0}
  \frac{\sqrt{s p_i^U}\,(\sqrt{1+u}\,u+u)}
       {|\operatorname{flsqrt}(\operatorname{flmul}(s,p_i^U))|}
\right)^2
\Biggr]
\sum_{j=1}^n
\left(
  \sum_{i:\,p_i^U>0}\frac{|U_{ij}|}{\sqrt{p_i^U}}
\right)^2 .
\end{aligned}
\]

\begin{theorem}[Equation (7), stored basis table and computed Gram]
Assume the hypotheses and sample-budget inequalities in the exact Bennett
baseline, and assume the length-\(s\) dot-product factor \(\gamma_s\) is valid.
Fix any storage mode \(\mu\in\mathcal M\).  Let
\(\Ghat_{\rm dot}^{(\mu)}(U)\) be the matrix whose entries are computed from
the exact leverage samples by using the stored table
\(\widehat U^{(\mu)}\), the concrete denominator table
\(\widehat d_i^U=\operatorname{flsqrt}(\operatorname{flmul}(s,p_i^U))\),
rounded row divisions, and rounded length-\(s\) Gram dot products.  Then
\[
  \Prob\!\left[
    \Ghat_{\rm dot}^{(\mu)}(U)-\I_n
      \preceq
      \bigl(\eps+T_U^{\rm fp}\bigr)\I_n
    \ \hbox{and}\
    \I_n-\Ghat_{\rm dot}^{(\mu)}(U)
      \preceq
      \bigl(\eps+T_U^{\rm fp}\bigr)\I_n
  \right]\ge 1-\delta .
\]
This is a final theorem for the three modeled stored-table paths
\(\operatorname{flmul}(U_{ij},1)\), \(\operatorname{fladd}(U_{ij},0)\), and
\(\operatorname{flsub}(U_{ij},0)\).  It is not a theorem about computing \(U\)
from a data matrix by QR, SVD, or another basis-generation routine.
\end{theorem}

\begin{remark}
In the small-roundoff regime with fixed \(U,m,n,s\), nonzero computed
positive-probability denominators, and
\(\widehat d_i^U/d_i^U\to1\), the denominator sum satisfies
\[
  \rho_U=\Theta\!\left(u\,|\{i:p_i^U>0\}|\right),
  \qquad
  \lambda_U=\Theta\!\left(u(1+|\{i:p_i^U>0\}|)\right).
\]
Thus the stored-basis radius has the compact order
\[
  T_U^{\rm fp}
  =
  \Theta\!\left(
    \bigl(\gamma_s+\lambda_U\bigr)
    \sum_{j=1}^n (B_j^U)^2
  \right),
\]
or, with the column magnitude expanded,
\[
  T_U^{\rm fp}
  =
  \Theta\!\left(
    \left(\gamma_s+u(1+|\{i:p_i^U>0\}|)\right)
    \sum_{j=1}^n
    \left(
      \sum_{i:\,p_i^U>0}\frac{|U_{ij}|}{\sqrt{p_i^U}}
    \right)^2
  \right).
\]
Since \(U^TU=I_n\) and \(n>1\), the displayed column factor is finite and
positive.  Therefore \(T_U^{\rm fp}\to0\) as \(u,\gamma_s\to0\) in this fixed
problem regime, so the stored-basis theorem is non-vacuous.
\end{remark}

\begin{remark}
The older vector-action consequence has probability
\[
  1-\frac{1}{s(\eps/n)^2}
\]
and the exact-denominator threshold
obtained by setting \(\rho_U=0\) and \(\eta_U=u\).  The Bennett
finite-Loewner results above are the current Algorithm 2 theorem surfaces for
the concrete computed-denominator path and for the three modeled stored-basis
paths.
\end{remark}

\section*{References}

\begin{thebibliography}{9}

\bibitem{DrineasMahoney2016RandNLA}
Petros Drineas and Michael W. Mahoney.
\newblock RandNLA: randomized numerical linear algebra.
\newblock \emph{Communications of the ACM}, 59(6):80--90, 2016.
\newblock \url{https://doi.org/10.1145/2842602}.

\bibitem{DrineasMahoneyMuthukrishnanSarlos2011}
Petros Drineas, Michael W. Mahoney, S. Muthukrishnan, and Tamas Sarlos.
\newblock Faster least squares approximation.
\newblock \emph{Numerische Mathematik}, 117:219--249, 2011.

\bibitem{Oliveira2010}
Roberto I. Oliveira.
\newblock Sums of random Hermitian matrices and an inequality by Rudelson.
\newblock arXiv:1004.3821, 2010.
\newblock \url{https://arxiv.org/abs/1004.3821}.

\bibitem{Tropp2012}
Joel A. Tropp.
\newblock User-friendly tail bounds for sums of random matrices.
\newblock \emph{Foundations of Computational Mathematics}, 12:389--434, 2012.
\newblock \url{https://doi.org/10.1007/s10208-011-9099-z}.

\bibitem{Higham2002}
Nicholas J. Higham.
\newblock \emph{Accuracy and Stability of Numerical Algorithms}.
\newblock Second edition, SIAM, 2002.
\newblock Standard source for the floating-point dot-product and
roundoff-error model used to state the deterministic perturbation budgets.

\end{thebibliography}

\end{document}
